\documentclass[11pt]{article}

\usepackage[margin=1in]{geometry}
\usepackage{graphicx}
\usepackage{float}
\usepackage{caption}
\usepackage{enumitem}
\usepackage{xcolor}
\usepackage[hidelinks]{hyperref}
\usepackage{parskip}

\captionsetup{font=small,labelfont=bf}
\setlist{itemsep=2pt,topsep=3pt}

\definecolor{accent}{HTML}{c0392b}
\newcommand{\fig}[3]{%
  \begin{figure}[H]
    \centering
    \includegraphics[width=#2\linewidth]{#1}
    \caption{#3}
  \end{figure}}

\title{\textbf{Inside the AI Conversation:\\ Sub-Category Drill-Downs on YouTube, 2026}}
\author{Pairie Koh \and Andrew B. Hall}
\date{June 3, 2026}

\begin{document}
\maketitle

\begin{abstract}
\noindent This memo opens up the top-level theme bars from our AI-sentiment
corpus on \textbf{YouTube} and asks what is \emph{inside} each one. The method
mirrors the companion TikTok memo: \textbf{Tier~1} uses fields already in the
corpus---the editorial \emph{register} of each video and the seed search keyword
that surfaced it---to characterise each theme without any new classification.
\textbf{Tier~2} runs a second, independent pass with Claude (Haiku~4.5, routed
through OpenRouter) that assigns every classified video to one theme-specific
sub-category. Everything here is \textbf{YouTube only} and restricted to videos
posted in \textbf{2026}, matching the main bar charts. Most of TikTok's lessons
survive: the deepfake fear is about \emph{scams}, not elections; the art backlash
is organised \#noAI activism; and the most genuine enthusiasm sits in AI
\emph{companionship}. But YouTube also adds two wrinkles. First, the
white-collar-tech story that was almost absent on TikTok is real on YouTube
(21\% of the jobs bar), reflecting CS YouTube's outsized share of platform
attention. Second, ``the public's environmental concern'' looks materially
different on YouTube than on TikTok---it is the electricity \emph{grid} (39\%),
not water, that leads.
\end{abstract}

\tableofcontents
\bigskip

\section{Method and scope}

The drill-downs sit on top of the corpus described in the main write-up:
AI-related YouTube videos collected via Bright Data, filtered and
theme-classified with Claude. We re-use the \emph{canonical} row filter from the
published bar charts (backlash $=$ \texttt{\_backlash\_final=YES} with a primary
topic; positive $=$ pollution dropped, \texttt{model\_stan}/\texttt{general\_hype}
secondaries promoted, \texttt{\_theme} valid), so the sub-category totals add up
to exactly the same bars readers already saw. In 2026, that filter yields
$n=2{,}397$ backlash videos and $n=4{,}704$ positive videos on YouTube.

\begin{itemize}
  \item \textbf{Tier 1 (no new classification).} The \emph{register} field
  (\texttt{genuine\_sentiment}, \texttt{creator\_promo},
  \texttt{creative\_showcase}, \texttt{news\_explainer}, \texttt{brand\_official})
  and the seed \texttt{\_search\_keyword} are already attached to each row.
  \item \textbf{Tier 2 (new sub-classification).} A second Claude pass assigns
  each video to one of a short, theme-specific menu of sub-categories. Hall's
  original TikTok pass used the Anthropic Message Batches API directly; the
  YouTube pass here is realtime-parallel through OpenRouter, but uses the same
  Haiku~4.5 model and an identical prompt/menu, so labels are comparable across
  platforms.
\end{itemize}

\paragraph{Caveats.} (i)~This is YouTube only; the companion TikTok memo
covers the other half. (ii)~The keyword breakdowns (Tier~1) reflect the seed
query, not content, and are skewed by how many seeds we used per theme---read
them as colour, not measurement. (iii)~In the 2026 cut, the model-fandom
($n=44$) and general-hype ($n=44$) positive sub-cells are too thin for headline
percentages. Existential risk ($n=126$) and breakthrough science ($n=181$) are
substantively larger on YouTube than on TikTok and can be reported, but should
still be treated as smaller-n than the headline themes.

\clearpage
\section{Tier 1: register and keyword structure}

\subsection{How organic is each positive theme?}

Splitting each positive theme by editorial register is again the most
consequential Tier~1 result, and again complicates any ``AI is well loved''
reading of the raw volume. On YouTube the high-volume themes are even more
promotional than on TikTok: career/productivity is 72\% creator promotion (course
funnels, tool demos, side-hustle pitches) and only 17\% genuine sentiment, and
creative tools are 52\% promotion plus another 31\% creative showcase, leaving
only 13\% expressing an actual opinion. As on TikTok, the genuine enthusiasm is
concentrated in the \emph{smaller} themes---AI companionship (41\% genuine) and
model fandom (39\% genuine)---though the absolute levels of genuineness are
lower than the equivalent TikTok cells (68\% and 85\%). Breakthrough science is
69\% news/explainer content, not public excitement; this is closer to
reposted-news framing than to original enthusiasm.

\fig{positive_register_split_2026_youtube.png}{1.0}{Register split of each
positive theme on YouTube. Green is genuine sentiment; orange and red are
promotional. The high-volume themes (top) are the most promotional; the most
genuine enthusiasm sits in AI companionship and model fandom, though at lower
absolute levels than the TikTok counterparts.}

\subsection{Seed-keyword breakdown per theme}

The keyword views give a fast, if rough, sense of what sits under each bar. On
the backlash side they confirm two YouTube-specific signals that matter for the
narrative: within the environment theme, ``AI data centers'' and ``AI
electricity'' jointly dominate the seed mix (consistent with the Tier~2
sub-distribution), and the art theme is again dominated by movement hashtags
(\#noAI, \#antiAI, ``stop AI art''). On the positive side, ``vibe coding'' and
``Cursor / Copilot'' lead career/productivity, ``Suno'' and ``Veo'' carry the
creative-tools theme, and the companion theme stretches across ``my AI
girlfriend,'' ``character.ai,'' and emotional-support framings.

\fig{drilldown_keywords_backlash_2026_youtube.png}{1.0}{Top seed search keywords
within each YouTube backlash theme (proxy sub-categories). Read as colour, not
measurement.}

\fig{drilldown_keywords_positive_2026_youtube.png}{1.0}{Top seed search keywords
within each YouTube positive theme. ``Vibe coding'' dominates career content;
``Suno'' and ``Veo'' lead creative tools; the companion theme spans
\texttt{character.ai}, ``AI girlfriend,'' and emotional-support framings.}

\clearpage
\section{Tier 2: independent sub-categories}

\subsection{Backlash}

The backlash drill-down delivers the same mass-versus-elite contrast as on
TikTok, and pushes some pieces of it harder. \textbf{Deepfakes} on YouTube are
even more concentrated as a scam fear: 51\% scam/fraud and 33\% general
misinformation, with political/election content at just 4.5\%---a category that
elite commentary expected to be the headline. \textbf{Jobs} anxiety is again
generalised (34\% ``all-jobs panic'' plus 25\% reaction to layoff news), but
unlike on TikTok, named \emph{white-collar tech} displacement is a real 21\% of
the bar, with service/blue-collar a rounding error (2\%). This is consistent
with CS YouTube's outsized share of platform attention---the white-collar story
shows up where the white-collar audience is. The \textbf{art} backlash is again
an organised movement: 53\% of it is \#noAI activism and identity content (a
shade lower than TikTok's 66\%, but still the modal sub-category by a wide
margin). Within \textbf{environment}, the ordering inverts relative to TikTok:
electricity/grid (39\%) leads, then data centers (30\%), with water usage a
distant third (15\%). \textbf{Existential risk}, finally large enough to read on
YouTube ($n=126$), is 67\% framed as \emph{loss of control}---``AI is out of
control,'' ``can't be stopped''---with the AGI/superintelligence framing only
11\% and explicit extinction-doom only 8\%; the public xrisk fear is closer to a
\emph{runaway technology} fear than to the rationalist AGI narrative.

\fig{drilldown_subtopics_backlash_2026_youtube.png}{1.0}{Sub-categories within
each backlash theme, YouTube 2026. Note the near-absence of election framing
inside deepfakes, the rebound of named white-collar inside jobs, and the
electricity-grid lead inside environment.}

\subsection{Positive}

On the positive side, the ``career win'' is again mostly a \emph{tooling demo}:
vibe coding (30\%) plus agent automation (20\%) plus ChatGPT-workflow tips
(11\%) together dwarf actually landing a job (3\%). Income / side-hustle pitches
are a striking 22\% of the bar on YouTube---larger than on TikTok---which is
consistent with the long-form-monetisation-funnel character of YouTube's
career-tutorial economy. AI companionship splits roughly 40/34/13/9 across
character roleplay, romantic, emotional support, and friendship---heavier on
character.ai-style roleplay than the TikTok cell. Model fandom is essentially a
four-way tie among Claude, ChatGPT, other named models, and generic amazement
(roughly 23\% each), with no Claude dominance equivalent to TikTok's 51\%.

\fig{drilldown_subtopics_positive_2026_youtube.png}{1.0}{Sub-categories within
each positive theme, YouTube 2026.}

\subsection{Selected themes in detail}

\fig{subtopic_jobs_2026_youtube.png}{0.82}{Jobs/displacement on YouTube.
Generalised panic and layoff news still dominate, but white-collar tech is a
substantial 21\% slice---visibly larger than its TikTok counterpart.}

\fig{subtopic_art_2026_youtube.png}{0.82}{Art/creative theft on YouTube. The
theme is again overwhelmingly organised \#noAI activism (53\%) rather than
individual job-loss complaints.}

\fig{subtopic_environment_2026_youtube.png}{0.82}{Environment/energy on YouTube.
Electricity/grid (39\%) leads, then data centers (30\%); water usage is a
distant third. The ordering is inverted relative to TikTok.}

\fig{subtopic_career_productivity_2026_youtube.png}{0.82}{Career/productivity on
YouTube. Vibe coding leads, side-hustle/income content is larger than on
TikTok, and reported employment outcomes remain a small slice.}

\fig{subtopic_companion_affective_2026_youtube.png}{0.82}{AI companion on
YouTube. Affective use is dominated by character roleplay and romantic framings
($\approx 74\%$ combined), with emotional support and friendship rounding out
the bar.}

\clearpage
\section{What this adds to the argument}

\begin{enumerate}
  \item \textbf{The deepfake fear is about being scammed, not about elections.}
  The elite frame (deepfakes endanger democracy) is essentially absent on
  YouTube too---in fact more starkly so than on TikTok (4.5\% vs 2\%, but
  against a much larger 51\% scam share).
  \item \textbf{The jobs panic is generalised, but with a real white-collar
  bulge on YouTube.} Unlike TikTok, named white-collar-tech displacement is a
  meaningful 21\% of the jobs bar---suggesting the platform on which the
  white-collar audience already lives is where the white-collar narrative finds
  purchase.
  \item \textbf{The environment story is about the electricity grid first.} On
  YouTube the resource ordering is electricity/grid $>$ data centers $>$
  water---the inverse of TikTok's water-leads pattern. The two platforms tell
  the public a different physical story about why AI is environmentally costly.
  \item \textbf{The art backlash is organised on YouTube too.} \#noAI activism
  is again the modal art sub-category by a wide margin, replicating the TikTok
  finding on a different audience.
  \item \textbf{Genuine enthusiasm is for companionship, not productivity.} Once
  promotion and showcase content are removed, the most heartfelt pro-AI content
  on YouTube is for AI companions and model loyalty---directionally identical to
  TikTok, just at lower absolute genuineness levels.
  \item \textbf{Existential risk is a ``runaway tech'' fear, not an AGI fear.}
  YouTube finally has enough xrisk content to read sub-cells: two-thirds is
  ``loss of control'' framing, with AGI/superintelligence at only 11\% and
  extinction at 8\%. The public xrisk language and the elite/rationalist xrisk
  language are not the same thing.
\end{enumerate}

\end{document}
